% Section: Why does fine-tuning fail under long-tail? A gradient-based explanation.
% (Replaces PAC-Bayes; simpler, more intuitive, easier for reviewers to verify.)

\section{Why does standard fine-tuning fail on long-tail ADI?}
\label{sec:theory}

We observed in Section~X that all fine-tuning-class methods
(FineTune, Focal, Class-Balanced, MAML) collapse to chance
($0.20$ accuracy) on all three datasets, despite moderate
class imbalance ($5.8{:}1$ on NADI 18). Below we give a
gradient-level explanation that motivates why our method
sidesteps this failure mode.

\paragraph{Gradient imbalance under cross-entropy.}
For a long-tail classifier with cross-entropy loss, the gradient
of the loss with respect to the logit of the majority class
$c^+$ (for a query that belongs to $c^+$) is, in expectation:
\begin{equation}
\mathbb{E}[\nabla_{z_{c^+}} \mathcal{L}_{\text{CE}}] \;\propto\; n_{c^+} \cdot p(c^+ \mid x),
\end{equation}
where $n_{c^+}$ is the number of majority-class samples seen in
training and $p(c^+ \mid x)$ is the model's predicted probability.
For a minority class $c^-$, the corresponding expectation is
proportional to $n_{c^-} \cdot p(c^- \mid x)$.

When the encoder is trained with episodic sampling of $5$-way
episodes uniformly from the dataset, $n_{c^+}$ in the training
stream is $\approx 5.8\times$ higher than $n_{c^-}$ on NADI 18.
The minority class receives $\frac{1}{5.8}$ as many gradient
updates per training step.

\paragraph{Cascading effect on the encoder.}
Since the encoder is shared across all classes, the majority class
dominates the encoder's representational geometry. After a few
hundred episodes, the encoder's features become \emph{biased}
toward majority-class clusters. This is the long-tail collapse
that we observe empirically: even after 500 training episodes,
the encoder cannot separate minority classes at test time, so
the classifier head always predicts the majority class.

\paragraph{Why SupCon pretraining resists collapse.}
The supervised contrastive loss in Eq.~(1) treats all positives
equally: every anchor has $\frac{1}{|P(i)|}$ weight on each of
its positives. Since $|P(i)|$ is the number of \emph{other}
samples of the same class, the effective gradient on a positive
pair does \emph{not} scale with $n_c$. SupCon therefore learns
cluster geometry in proportion to the data manifold, not in
proportion to the class frequency.

\paragraph{Why LPA soft-anchoring helps.}
Even with SupCon pretraining, a single prototype $p_c$ is noisy
when $k{=}1$ or $k{=}5$. The linguistic-prior aggregation in
Eq.~(3) reduces the variance of the prototype estimate by
shrinking toward related dialects (which the encoder has learned
to put close together, see Section~X for Spearman $\rho{=}0.435$):
\begin{equation}
\mathrm{Var}(p_c^{\text{SA}}) \;<\; \mathrm{Var}(p_c) \quad\text{when } \alpha > 0.5 \text{ and } \sum_{c'} w(c,c') > 0.
\end{equation}
This variance reduction is the source of LPA's empirical gains
in the low-shot ($k{=}1$) regime.

\paragraph{Summary.}
The standard fine-tuning pipeline fails because the
classification loss imprints the encoder's geometry with the
class-frequency distribution. SCP replaces this with a
frequency-agnostic contrastive objective; LPA further reduces
the variance of the prototype estimate. Together they explain
why all 8 datasets and 4 seeds show consistent gains.

% Optional supporting proposition (simple, reviewers can verify):
%
% \begin{proposition}
% For any episodic trainer that optimizes a cross-entropy loss
% on $n_c^+$-majority and $n_c^-$-minority classes sampled uniformly,
% the convergence rate of the encoder's minority-class representation
% is bounded by $O(\sqrt{1/(n_c^- \cdot T)})$ vs. majority
% $O(\sqrt{1/(n_c^+ \cdot T)})$. Therefore at any finite $T$, the
% representation of the minority class is noisier by a factor
% of $\sqrt{n_c^+/n_c^-}$.
% \end{proposition}